\documentclass[11pt,letterpaper]{article}

\usepackage[utf8]{inputenc}
\usepackage[T1]{fontenc}
\usepackage[spanish,es-noquoting]{babel}
\usepackage{lmodern}
\usepackage[margin=2.5cm]{geometry}

\usepackage{microtype}
\usepackage{graphicx}
\usepackage{array}
\usepackage{tabularx}
\usepackage{booktabs}
\usepackage{tikz}
\usepackage{enumitem}
\usepackage{listings}
\usepackage[hidelinks]{hyperref}

\usetikzlibrary{arrows.meta,positioning,calc}

\setlength{\parindent}{0pt}
\setlength{\parskip}{6pt}
\setlength{\emergencystretch}{2em}

\setlist{nosep,leftmargin=1.4em}

\pagestyle{plain}

\lstset{
    basicstyle=\ttfamily\small,
    breaklines=true,
    columns=fullflexible,
    keepspaces=true,
    showstringspaces=false
}

\newcommand{\archivo}[1]{\texttt{\detokenize{#1}}}

\newcommand{\crc}[3]{%
    \fbox{%
        \begin{minipage}[t]{.445\textwidth}
            \small
            \textbf{#1}

            \textbf{Responsabilidades}

            #2

            \textbf{Colaboradores}

            #3
        \end{minipage}%
    }%
}

\tikzset{
    modulo/.style={
        draw,
        align=center,
        font=\small,
        minimum height=1cm,
        text width=4cm
    },
    flecha/.style={
        -{Stealth[length=2mm]}
    },
    etiqueta/.style={
        font=\scriptsize,
        fill=white,
        inner sep=2pt
    }
}

\begin{document}

\begin{center}
    \textbf{Universidad Nacional Autónoma de México}\\
    Facultad de Ciencias

    \textbf{Reporte de Modelado y Programación}\\
    Base de datos musical con interfaz gráfica en C\# y Avalonia UI (RolasApp)

    Guillermo Ricardo Gómez Gómez
\end{center}

\section{Resumen}

Desarrollé \textbf{RolasApp}, una aplicación de escritorio para construir y
administrar una base de datos musical a partir de directorios MP3 con
etiquetas ID3v2.4, utilizando C\# (\archivo{net10.0}), Avalonia UI, SQLite
y TagLibSharp. El proyecto tuvo como propósito integrar el modelado
relacional con una interfaz gráfica multiplataforma siguiendo la
arquitectura Modelo--Vista--Controlador (MVC). La principal dificultad
de diseño consistió en equilibrar el modelado inicial pensado en cuaderno
con una estructura práctica y mantenible, evitando capas redundantes y
garantizando que ninguna modificación alterara los directorios MP3 físicos
en disco. Para resolverlo, separé la persistencia en tres objetos de
acceso a datos (DAO), encapsulé el escaneo recursivo y la lectura de
metadatos en un minero dedicado (\archivo{Steve}) y dividí la interfaz en
dos ventanas especializadas. Como resultado, la aplicación mina colecciones
completas en segundo plano, muestra portadas y fotografías de artistas
incrustadas, permite filtrar pistas y gestiona tanto solistas como bandas
musicales con sus respectivos integrantes.

\section{Introducción}

El objetivo del segundo proyecto de la materia fue desarrollar un
sistema capaz de poblar automáticamente una base de datos relacional
a partir del recorrido de un directorio con música en formato MP3 y
ofrecer una interfaz gráfica para consultar y editar dicha información.
A diferencia del primer proyecto, donde trabajé con sockets en C y un
cliente de consola en C\#, en esta ocasión me enfoqué por completo en el
ecosistema de .NET sobre Linux, incorporando bibliotecas externas para
el acceso a bases de datos embebidas, la lectura de cabeceras binarias
ID3v2.4 y la construcción de ventanas gráficas nativas.

Aunque ya contaba con experiencia previa trabajando con C\# en consola,
desarrollar una aplicación con interfaz gráfica en Avalonia UI implicó
aprender el enlace de datos (\textit{data binding}), la separación entre
la definición declarativa en XAML y la lógica de eventos en C\#, así como
la ejecución de tareas pesadas en hilos secundarios para evitar que la
ventana dejara de responder durante el escaneo de directorios grandes.
Asimismo, tuve que respetar los lineamientos del estándar XDG de Linux
para ubicar la base de datos en \archivo{~/.local/share/RolasApp/rolas.db}
y resolver la ruta predeterminada de música del usuario sin depender de
nombres de carpetas fijos.

La documentación técnica fue un apoyo constante durante todo el proceso.
Consulté la documentación oficial de Avalonia UI para configurar estilos
de listas, selectores nativos de carpetas mediante \archivo{StorageProvider}
y ventanas modales; la guía de \archivo{Microsoft.Data.Sqlite} para
manejar conexiones, consultas parametrizadas y transacciones ACID; los
ejemplos de TagLibSharp para extraer etiquetas ID3v2.4 y marcos gráficos
\archivo{APIC}; y la especificación oficial de ID3v2.4 para verificar el
significado de cada identificador de trama \cite{avalonia,sqlite-ms,taglib,id3v24}.

Al comenzar el desarrollo, cometí el error de implementar escritura de
regreso sobre las etiquetas de los directorios MP3 cada vez que el usuario
guardaba un cambio en la interfaz. Al revisar con detenimiento los
requerimientos del proyecto, me di cuenta que los directorios MP3 debían
permanecer estrictamente intactos y que toda edición debía ocurrir
únicamente dentro de SQLite. Arreglar esto, junto con acomodar mejor
las ventanas y completar el manejo de personas y bandas, fue lo que hizo
que el proyecto pasara de ser un primer avance funcional a cumplir bien
con lo que pedía la entrega.

\section{Análisis y diseño}

\subsection{Del planteamiento inicial a una arquitectura MVC}

Antes de escribir código, tracé en cuaderno un diagrama general de
clases y paquetes dividido en tres columnas: \archivo{Vista},
\archivo{Modelo} y \archivo{Controlador}. En ese diseño inicial había
contemplado cuatro elementos que después decidí replantear: una carpeta
intermedia \archivo{Vista/Contratos} con interfaces para desacoplar
Avalonia de los controladores, un subdirectorio \archivo{Modelo/Compilador}
destinado a realizar análisis léxico y semántico de consultas para
traducirlas a SQL, y un tercer controlador llamado \archivo{C. Reproductor}
para la reproducción de audio.

Cuando comencé a trasladar ese diagrama al repositorio de .NET, creé las
carpetas vacías \archivo{Vista/Contratos} y \archivo{Modelo/Compilador}
mientras avanzaba primero con las entidades, la conexión a SQLite y el
minero. Pronto noté que mantener carpetas vacías y diseñar un analizador
léxico propio para búsquedas sencillas y sobre todo interfaces de contrato
para solo dos ventanas añadía complejidad al proyecto y no podria concluirlo
en tiempo y forma. Además, tanto el compilador de consultas como el reproductor
eran componentes opcionales en la especificación.

\begin{figure}[htbp]
\centering
\begin{tikzpicture}
    \node[modulo,text width=4.2cm] (vista) at (-4.9,0) {
        \textbf{Vista (UI)}\\
        Ventana Principal\\
        Editar Grupo / Miembros\\
        \textit{Contratos (interfaces)}
    };

    \node[modulo,text width=4.4cm] (ctrl) at (0,0) {
        \textbf{Controlador}\\
        Cont. Principal\\
        C. Intérpretes\\
        \textit{C. Reproductor (opcional)}
    };

    \node[modulo,text width=4.4cm] (mod) at (4.9,0) {
        \textbf{Modelo}\\
        Minero (\archivo{Steve})\\
        Entidades y DAOs\\
        \textit{Compilador SQL (opcional)}
    };

    \draw[flecha,<->] (vista)--(ctrl);
    \draw[flecha,<->] (ctrl)--(mod);
\end{tikzpicture}
\caption{Planteamiento inicial de la arquitectura antes de depurar módulos opcionales y contratos.}
\end{figure}

Decidí entonces simplificar el diseño, eliminé las carpetas vacías,
sustituí el compilador léxico por un motor de filtrado en memoria
dentro de \archivo{ControladorRolas} capaz de interpretar prefijos
como \archivo{artista:}, \archivo{album:} o \archivo{año:} sin que
el usuario escriba SQL, y concentré el esfuerzo en las dos partes
centrales del sistema: el catálogo de pistas y la administración
de intérpretes, individuos y bandas.

\subsection{Organización final del proyecto}

Tras la refactorización y la reorganización de las ventanas dentro del
espacio de nombres \archivo{RolasApp.Vista}, la estructura final del
proyecto quedó dividida de forma simétrica en tres capas bien
delimitadas. La capa de \archivo{Vista} contiene exclusivamente las
definiciones XAML y el código subyacente de
\archivo{MainWindow} y \archivo{VentanaInterpretes}. La capa de
\archivo{Controlador} agrupa a \archivo{ControladorRolas} y
\archivo{ControladorInterpretes}, encargados de coordinar las
validaciones, las búsquedas y las llamadas a la persistencia. Finalmente,
la capa de \archivo{Modelo} reúne las entidades del dominio, los tres
DAO sobre \archivo{ConexionBaseDatos} y el minero \archivo{Steve}.

\begin{figure}[htbp]
\centering
\begin{tikzpicture}
    \node[modulo,text width=5.2cm] (mw) at (-3.4,0) {
        \textbf{«vista» MainWindow}\\
        Catálogo, buscador y edición en BD
    };

    \node[modulo,text width=5.2cm] (vi) at (3.4,0) {
        \textbf{«vista» VentanaInterpretes}\\
        Personas, bandas y miembros
    };

    \node[modulo,text width=5.2cm] (cr) at (-3.4,-2.2) {
        \textbf{«controlador» ControladorRolas}\\
        Minado asíncrono, filtrado y CRUD
    };

    \node[modulo,text width=5.2cm] (ci) at (3.4,-2.2) {
        \textbf{«controlador» ControladorInterpretes}\\
        Clasificación y relación \archivo{in_group}
    };

    \node[modulo,text width=4.0cm] (st) at (-4.8,-4.5) {
        \textbf{«minero» Steve}\\
        TagLibSharp (ID3v2.4 / APIC)
    };

    \node[modulo,text width=4.2cm] (daos) at (0,-4.5) {
        \textbf{«persistencia» DAOs}\\
        \archivo{DaoRolas}, \archivo{DaoAlbumes},\\
        \archivo{DaoInterpretes}
    };

    \node[modulo,text width=4.0cm] (ent) at (4.8,-4.5) {
        \textbf{«modelo» Entidades}\\
        \archivo{Rola}, \archivo{Album}, \archivo{Interprete},\\
        \archivo{Persona}, \archivo{Grupo}
    };

    \node[modulo,text width=6.4cm] (bd) at (0,-6.7) {
        \textbf{«singleton» ConexionBaseDatos / SQLite}\\
        \archivo{~/.local/share/RolasApp/rolas.db}
    };

    \draw[flecha] (mw)--(cr);
    \draw[flecha] (mw)--node[etiqueta,above]{abre diálogo}(vi);
    \draw[flecha] (vi)--(ci);
    \draw[flecha] (cr)--(st);
    \draw[flecha] (cr)--(daos);
    \draw[flecha] (ci)--(daos);
    \draw[flecha] (st)--(daos);
    \draw[flecha] (daos)--(ent);
    \draw[flecha] (daos)--(bd);
\end{tikzpicture}
\caption{Organización definitiva de las capas MVC en \textbf{RolasApp}.}
\end{figure}

En la siguiente tabla se resumen las principales decisiones de diseño que
tomé durante el desarrollo y la justificación de cada una:

\begin{center}
\small
\begin{tabularx}{\textwidth}{@{}p{4.0cm}X@{}}
    \toprule
    Decisión & Motivo \\
    \midrule

    Eliminación de contratos XAML &
    Para un proyecto de dos ventanas con entrega acotada, vincular las
    vistas directamente con los controladores evitó duplicar firmas de
    métodos sin perder la separación de capas MVC. \\

    Búsqueda por prefijos en memoria &
    Permite combinar filtros como \archivo{artista:Queen año:1975} o
    texto libre en tiempo real sin necesidad de construir un analizador
    léxico ni exponer SQL al usuario. \\

    Tabla de catálogo estilo Finder &
    El uso de filas alternadas (\textit{zebra striping}) y columnas
    alineadas mejora la legibilidad cuando la biblioteca contiene cientos
    de pistas minadas. \\

    Edición exclusiva en SQLite &
    Garantiza el cumplimiento estricto de la especificación, manteniendo
    los directorios MP3 originales en modo de solo lectura. \\

    Ubicación según estándar XDG &
    Almacenar \archivo{rolas.db} en \archivo{LocalApplicationData} y
    consultar \archivo{SpecialFolder.MyMusic} integra la aplicación de
    forma limpia en entornos Linux. \\

    \bottomrule
\end{tabularx}
\end{center}

\subsection{Modelo relacional y esquema de la base de datos}

La base de datos sigue exactamente el esquema relacional de siete tablas
establecido en la especificación (\archivo{rolas.sql}). La tabla
\archivo{types} actúa como catálogo fijo para distinguir si un
intérprete registrado en \archivo{performers} corresponde a una persona
individual (\archivo{0}), a una agrupación musical (\archivo{1}) o a un
intérprete aún sin clasificar (\archivo{2}, valor asignado por defecto
cuando \archivo{Steve} descubre un artista nuevo durante el minado).

Por su parte, las tablas \archivo{persons} y \archivo{groups} almacenan
los atributos propios de cada tipo de intérprete: nombre artístico,
nombre real y fechas de nacimiento o fallecimiento para los individuos; y
nombre de la banda junto con sus fechas de inicio y separación para los
grupos. La relación muchos a muchos entre músicos y agrupaciones se
materializa mediante la tabla intermedia \archivo{in_group}, cuya llave
primaria compuesta \archivo{(id_person, id_group)} impide registrar dos
veces al mismo integrante dentro de una banda.

\begin{figure}[htbp]
\centering
\begin{tikzpicture}
    \node[modulo,text width=3.9cm] (types) at (-5.2,0) {
        \textbf{types}\\
        \underline{id\_type}, description
    };

    \node[modulo,text width=4.3cm] (perf) at (0,0) {
        \textbf{performers}\\
        \underline{id\_performer}, id\_type,\\
        name
    };

    \node[modulo,text width=4.1cm] (alb) at (5.2,0) {
        \textbf{albums}\\
        \underline{id\_album}, path,\\
        name, year
    };

    \node[modulo,text width=4.3cm] (pers) at (-5.2,-2.5) {
        \textbf{persons}\\
        \underline{id\_person}, stage\_name,\\
        real\_name, birth\_date,\\
        death\_date
    };

    \node[modulo,text width=4.1cm] (grp) at (0,-2.5) {
        \textbf{groups}\\
        \underline{id\_group}, name,\\
        start\_date, end\_date
    };

    \node[modulo,text width=4.5cm] (rolas) at (5.2,-2.5) {
        \textbf{rolas}\\
        \underline{id\_rola}, id\_performer,\\
        id\_album, path, title,\\
        track, year, genre
    };

    \node[modulo,text width=4.3cm] (ingrp) at (-2.6,-5.0) {
        \textbf{in\_group}\\
        \underline{id\_person, id\_group}
    };

    \draw[flecha] (perf.west)--node[etiqueta,above]{FK}(types.east);
    \draw[flecha] (rolas.north west)--node[etiqueta,above]{FK}(perf.south east);
    \draw[flecha] (rolas.north)--node[etiqueta,right]{FK}(alb.south);
    \draw[flecha] (ingrp.north west)--node[etiqueta,left]{FK}(pers.south);
    \draw[flecha] (ingrp.north east)--node[etiqueta,right]{FK}(grp.south);
\end{tikzpicture}
\caption{Esquema relacional de siete tablas implementado en SQLite.}
\end{figure}

Finalmente, \archivo{albums} guarda la ruta del directorio contenedor en
el sistema de archivos, el título del disco y su año de publicación,
mientras que \archivo{rolas} concentra los metadatos de cada pista
individual junto con la ruta absoluta hacia el directorio MP3 y las llaves
foráneas hacia \archivo{performers} y \archivo{albums}.

\subsection{Tarjetas CRC como revisión del diseño}

Para verificar que cada clase del sistema tuviera una responsabilidad
acotada y evitar que los controladores o las ventanas acumularan lógica
de acceso a datos, elaboré tarjetas CRC sobre los ocho componentes
principales de la aplicación:

\noindent
\crc{Steve (\archivo{Steve.cs})}
    {Recorrer directorios recursivamente, leer etiquetas ID3v2.4 y
    marcos \archivo{APIC}, y poblar la base de datos dentro de una
    transacción.}
    {TagLibSharp, \archivo{ConexionBaseDatos}, \archivo{DaoInterpretes},
    \archivo{DaoAlbumes}, \archivo{DaoRolas} y \archivo{ConstructorRola}.}
\hfill
\crc{ConexionBaseDatos}
    {Mantener una única conexión activa con SQLite (Singleton), activar
    llaves foráneas, crear tablas y coordinar transacciones ACID.}
    {SqliteConnection, SqliteTransaction y los tres DAOs.}

\vspace{7pt}
\noindent
\crc{DaoRolas y DaoAlbumes}
    {Insertar, consultar y actualizar pistas y álbumes evitando
    duplicados por ruta en disco.}
    {\archivo{ConexionBaseDatos}, \archivo{Rola} y \archivo{Album}.}
\hfill
\crc{DaoInterpretes}
    {Registrar intérpretes, clasificarlos como Persona o Grupo,
    administrar individuos sin pistas y gestionar \archivo{in_group}.}
    {\archivo{ConexionBaseDatos}, \archivo{Interprete}, \archivo{Persona}
    y \archivo{Grupo}.}

\vspace{7pt}
\noindent
\crc{ControladorRolas}
    {Ejecutar el minado en hilo secundario con reporte de progreso,
    filtrar rolas por prefijos y actualizar pistas y álbumes en SQLite.}
    {\archivo{Steve}, \archivo{DaoRolas}, \archivo{DaoAlbumes},
    \archivo{DaoInterpretes} y \archivo{MainWindow}.}
\hfill
\crc{ControladorInterpretes}
    {Coordinar el cambio de tipo de intérprete, el registro de
    individuos independientes y la asignación de integrantes a bandas.}
    {\archivo{DaoInterpretes} y \archivo{VentanaInterpretes}.}

\vspace{7pt}
\noindent
\crc{MainWindow}
    {Mostrar la tabla de rolas, filtrar en tiempo real, abrir el
    selector nativo de carpetas y desplegar portada y foto del artista.}
    {\archivo{ControladorRolas}, \archivo{VentanaInterpretes} y
    Avalonia UI.}
\hfill
\crc{VentanaInterpretes}
    {Permitir clasificar intérpretes, editar datos de personas y grupos,
    dar de alta músicos sin pistas y gestionar miembros de bandas.}
    {\archivo{ControladorInterpretes} y Avalonia UI.}

\subsection{Flujo de minado asíncrono e integridad de los directorios MP3}

Uno de los retos técnicos de las aplicaciones de escritorio es evitar que
una operación intensiva de entrada/salida congele la interfaz gráfica.
Si el minado de varios cientos de directorios MP3 se ejecutara en el hilo
principal de Avalonia, la ventana dejaría de repintarse hasta concluir
todo el recorrido. Para resolverlo, encapsulé la llamada a
\archivo{Steve.Minar()} dentro de \archivo{Task.Run()} en
\archivo{ControladorRolas.MinarAsync()} y utilicé \archivo{Progress<int>}
para actualizar la barra de avance en la barra inferior de
\archivo{MainWindow} de manera segura entre hilos.

Asimismo, realizar una escritura individual en SQLite por cada directorio MP3
degradaba el rendimiento debido al costo de sincronización en disco por
cada sentencia \archivo{INSERT}. Al envolver el ciclo completo de
\archivo{ProcesarMp3()} dentro de \archivo{IniciarTransaccion()} y
\archivo{ConfirmarTransaccion()}, todas las inserciones se confirman en
una sola operación atómica, reduciendo drásticamente el tiempo de
escaneo y asegurando que, ante cualquier fallo inesperado, la base de
datos pueda revertirse sin quedar a medias.

\begin{figure}[htbp]
\centering
\begin{tikzpicture}[x=1cm,y=1cm]
\foreach \x/\nombre in {
    0/MainWindow,
    3.6/ControladorRolas,
    7.2/Steve,
    10.8/DAOs,
    14.2/SQLite
}{
    \node[
        font=\small\bfseries,
        align=center
    ] at (\x,0) {\nombre};

    \draw[dashed,gray] (\x,-.3)--(\x,-4.5);
}

\draw[flecha]
    (0,-.7)--node[etiqueta,above]
    {MinarAsync(ruta)}(3.6,-.7);

\draw[flecha]
    (3.6,-1.3)--node[etiqueta,above]
    {Task.Run(Minar)}(7.2,-1.3);

\draw[flecha]
    (7.2,-1.9)--node[etiqueta,above]
    {BEGIN TRANSACTION}(14.2,-1.9);

\draw[flecha]
    (7.2,-2.6)--node[etiqueta,above]
    {ObtenerOCrear / Guardar}(10.8,-2.6);

\draw[flecha]
    (7.2,-3.2)--node[etiqueta,above]
    {Progress(porcentaje)}(0,-3.2);

\draw[flecha]
    (7.2,-3.9)--node[etiqueta,above]
    {COMMIT}(14.2,-3.9);
\end{tikzpicture}
\caption{Secuencia del minado asíncrono con reporte de progreso y transacción única en SQLite.}
\end{figure}

Respecto a la integridad de los directorios originales, \archivo{Steve} abre
cada directorio MP3 únicamente en modo lectura mediante
\archivo{TagLib.File.Create()} para extraer sus etiquetas e imágenes
\archivo{APIC}. Cuando el usuario edita el título, el intérprete, el
álbum o el año desde \archivo{MainWindow}, \archivo{ControladorRolas}
actualiza exclusivamente los registros correspondientes en las tablas
\archivo{rolas}, \archivo{albums} y \archivo{performers}, dejando intacto
el directorio físico en el sistema de archivos.

\section{Implementación}

\subsection{Capa de persistencia, transacciones y patrones de diseño}

En la capa de persistencia implementé \archivo{ConexionBaseDatos}
aplicando el patrón \textit{Singleton} con inicialización perezosa
(\archivo{Lazy<ConexionBaseDatos>}), lo que garantiza que todos los DAOs
compartan exactamente la misma conexión abierta y participen en la misma
transacción activa cuando \archivo{Steve} realiza un escaneo masivo. Al
abrir la conexión, el constructor ejecuta \archivo{PRAGMA foreign_keys = ON;}
para que SQLite haga cumplir la integridad referencial definida en el
esquema y verifica en \archivo{sqlite_master} si las tablas ya existen;
en caso contrario, ejecuta el script de creación automáticamente.

Para evitar duplicados al minar varias veces una misma carpeta,
\archivo{DaoRolas.GuardarRola()} consulta primero si ya existe una fila
con la misma ruta absoluta (\archivo{path}); si existe, actualiza sus
campos en lugar de insertar un registro repetido. De manera análoga,
\archivo{DaoAlbumes.ObtenerOCrearAlbum()} identifica cada álbum por el
par \archivo{(path, name)}, infiriendo como ruta del álbum el directorio
padre que contiene los directorios MP3 de esa pista.

En la construcción de entidades apliqué el patrón \textit{Builder} a
través de la clase \archivo{ConstructorRola} dentro de \archivo{Rola.cs}.
Este constructor fluido valida que ninguna pista se instancie sin una
ruta de directorio válida y normaliza los campos vacíos o nulos al valor por
omisión \archivo{"Unknown"}, manteniendo el código de \archivo{Steve}
limpio y expresivo.

\subsection{Minado de etiquetas ID3v2.4 e imágenes APIC con Steve}

El recorrido de directorios se realiza mediante
\archivo{Directory.EnumerateFiles()} configurado con
\archivo{EnumerationOptions}, activando la búsqueda recursiva, la
insensibilidad a mayúsculas y minúsculas para localizar tanto \archivo{.mp3}
como \archivo{.MP3}, y la opción \archivo{IgnoreInaccessible = true} para
evitar que el escaneo se interrumpa si encuentra subcarpetas protegidas
por permisos del sistema operativo.

Para cada directorio encontrado, \archivo{Steve} lee las seis etiquetas
ID3v2.4 solicitadas en la especificación mediante TagLibSharp:
\archivo{FirstPerformer} (\archivo{TPE1}), \archivo{Album} (\archivo{TALB}),
\archivo{Title} (\archivo{TIT2}), \archivo{FirstGenre} (\archivo{TCON}),
\archivo{Track} (\archivo{TRCK}) y \archivo{Year} (\archivo{TDRC})
\cite{taglib,id3v24}. Si alguna cabecera está dañada o el directorio no es
un MP3 válido, los bloques \archivo{catch} específicos para
\archivo{CorruptFileException} y \archivo{UnsupportedFormatException}
descartan esa pista sin abortar el resto del lote.

Además de los metadatos textuales, incorporé en \archivo{Steve} dos
métodos de lectura binaria sobre el arreglo \archivo{mp3.Tag.Pictures}:
\archivo{ObtenerPortada()}, que prioriza el marco \archivo{APIC} de tipo
\archivo{FrontCover}, y \archivo{ObtenerFotoArtista()}, que busca marcos
de tipo \archivo{Artist} o \archivo{LeadArtist}. Los bytes extraídos se
convierten en memoria a objetos \archivo{Bitmap} de Avalonia cada vez que
el usuario selecciona una pista en la tabla.

\subsection{Interfaz gráfica en Avalonia UI y gestión de intérpretes y bandas}

La ventana principal (\archivo{MainWindow}) está organizada en tres zonas:
una barra superior con la caja de búsqueda y los botones de acción, un
panel central dividido entre la tabla de rolas y el editor lateral, y
una barra inferior con mensajes de estado y una barra de progreso. En la
barra superior integré un botón para minar directamente el directorio
musical del usuario (\archivo{Steve.ObtenerRutaMusica()}) y otro botón
(\archivo{Elegir Directorio...}) que abre el diálogo nativo del escritorio
mediante \archivo{StorageProvider.OpenFolderPickerAsync()} para analizar
cualquier carpeta del sistema.

Para resolver el requerimiento de clasificar intérpretes, editar sus
detalles, dar de alta individuos sin pistas y gestionar miembros de
bandas, implementé \archivo{VentanaInterpretes} respaldada por
\archivo{ControladorInterpretes}. Esta ventana secundaria presenta a la
izquierda el listado de todos los intérpretes registrados en
\archivo{performers}; al seleccionar uno, el usuario puede definir
mediante un \archivo{ComboBox} si se trata de una \textbf{Persona} o un
\textbf{Grupo}. Al elegir \textbf{Persona}, se habilitan los campos para
editar su nombre artístico, nombre real y fechas de nacimiento y
fallecimiento en \archivo{persons}. Al elegir \textbf{Grupo}, se despliegan
los campos de la banda en \archivo{groups} junto con un selector de
personas registradas para agregar o retirar integrantes en
\archivo{in_group}. Finalmente, un tercer bloque en la misma ventana
permite dar de alta individuos directamente en \archivo{persons} aunque no
interpreten ninguna pista en el catálogo.

En cuanto a la automatización de compilación, organicé todas las tareas en
un único \archivo{Makefile} donde \archivo{make compilar} y
\archivo{make ejecutar} operan exclusivamente sobre la solución de .NET,
mientras que \archivo{make ejecutar reporte} compila este documento con
\archivo{pdflatex} dentro de \archivo{docs/} y abre el PDF resultante de
forma independiente.

\section{Conclusiones}

A través del desarrollo de \textbf{RolasApp} logré integrar el diseño de
una base de datos relacional normalizada con una interfaz gráfica de
escritorio funcional, aplicando patrones como \textit{Singleton},
\textit{Builder} y DAO dentro de una arquitectura MVC. Comprobar cómo los
datos extraídos automáticamente de directorios binarios MP3 poblaban las
tablas de SQLite y se reflejaban al instante en la interfaz confirmó la
importancia de separar claramente la lógica de acceso a datos de la
presentación visual.

\subsection{Dificultades y correcciones durante el desarrollo}

Durante la implementación surgieron varios errores concretos que me
obligaron a revisar tanto el diseño como los detalles técnicos del código:

\begin{itemize}
    \item \textbf{Escritura accidental sobre los directorios MP3:} En una
    primera versión implementé un método \archivo{ActualizarEtiquetasMp3}
    en \archivo{Steve} que invocaba \archivo{mp3.Save()} al editar una
    pista. Al releer la especificación y notar que los directorios MP3 no
    debían modificarse bajo ninguna circunstancia, eliminé por completo
    dicho método y restringí las actualizaciones a SQLite.
    \item \textbf{Duplicación de registros al editar nombres en los DAOs:}
    En \archivo{DaoInterpretes}, los métodos \archivo{GuardarPersona} y
    \archivo{GuardarGrupo} buscaban inicialmente registros previos solo
    por nombre. Cuando el usuario corregía el nombre artístico de una
    persona o el nombre de una banda existente, el DAO insertaba un
    registro nuevo en lugar de actualizar el actual. Lo corregí
    priorizando la búsqueda por llave primaria (\archivo{IdPersona} e
    \archivo{IdGrupo}) antes de consultar por nombre.
    \item \textbf{Errores de sintaxis en XAML y enlace de eventos:} Al
    ajustar los estilos de \archivo{ListBox} en \archivo{MainWindow.axaml},
    olvidé cerrar las etiquetas \archivo{</Grid>} y \archivo{</Border>}
    del panel izquierdo y referencié manejadores de eventos que aún no
    existían en el directorio \archivo{.axaml.cs}, provocando errores de
    compilación \archivo{AVLN1001} y \archivo{AVLN3000}.
    \item \textbf{Sincronización de espacios de nombres al reorganizar las vistas:}
    Al mover \archivo{MainWindow} hacia la carpeta \archivo{Vista/}, fue
    necesario actualizar el atributo \archivo{x:Class} en los directorios
    XAML e importar \archivo{RolasApp.Vista} en \archivo{App.axaml.cs}
    para que el punto de entrada de Avalonia localizara la ventana
    principal.
\end{itemize}

\subsection{Cosas a mejorar}

Aunque la versión entregada cumple con todos los requerimientos
obligatorios y ofrece una experiencia fluida, identifiqué tres mejoras
que sería valioso incorporar en una iteración posterior:

\begin{itemize}
    \item \textbf{Reproductor de audio integrado:} Incorporar el
    controlador \archivo{C. Reproductor} planteado en el diseño original
    utilizando una biblioteca multiplataforma de audio para escuchar las
    pistas directamente desde la ventana principal sin depender de un
    reproductor externo del sistema.
    \item \textbf{Consultas booleanas avanzadas:} Extender el buscador
    actual para soportar operadores lógicos (\archivo{AND}, \archivo{OR},
    rangos de años como \archivo{año:1970..1980}) mediante un pequeño
    analizador sintáctico, retomando la idea original del módulo
    \archivo{Compilador} pero de forma acotada.
    \item \textbf{Pruebas unitarias:} Añadir un proyecto de
    pruebas con xUnit dentro de la solución para verificar de manera
    automática las transacciones de los DAOs, el filtrado por prefijos y
    la extracción de etiquetas de \archivo{Steve} ante directorios corruptos o
    con metadatos incompletos.
\end{itemize}

\begingroup
\small

\begin{thebibliography}{9}

\bibitem{avalonia}
Avalonia UI Team.
\textit{Avalonia UI Documentation: Controls, Styling and StorageProvider}.
\url{https://docs.avaloniaui.net/}.

\bibitem{sqlite-ms}
Microsoft.
\textit{Información general de Microsoft.Data.Sqlite y transacciones}.
Microsoft Learn.
\url{https://learn.microsoft.com/es-es/dotnet/standard/data/sqlite/}.

\bibitem{taglib}
Mono Project / TagLibSharp Contributors.
\textit{TagLib\# (TagLibSharp) Library Documentation}.
Repositorio oficial en GitHub.
\url{https://github.com/mono/taglib-sharp}.

\bibitem{id3v24}
Martin Nilsson.
\textit{ID3 tag version 2.4.0 - Native Frames}.
Especificación formal de marcos ID3v2.4 (\archivo{TPE1}, \archivo{TALB}, \archivo{TIT2}, \archivo{TCON}, \archivo{TRCK}, \archivo{TDRC}, \archivo{APIC}).
\url{https://id3.org/id3v2.4.0-frames}.

\bibitem{xdg}
Waldo Bastian, Ryan Lortie y Lennart Poettering.
\textit{XDG Base Directory Specification}.
freedesktop.org.
\url{https://specifications.freedesktop.org/basedir-spec/latest/}.

\bibitem{notas-myp}
Kimberly Becerril Fuentes y Canek Peláez Valdés.
\textit{Notas para el curso de Modelado y Programación}.
Facultad de Ciencias, Universidad Nacional Autónoma de México, 2019.

\end{thebibliography}

\endgroup
\end{document}